\section{Constructing the Barycentric Interaction Field}
\label{sec:interaction-operator}

The preceding section shows how an admissible peer matrix enters the
exposure-adjustment model.
We now ask what should populate that matrix when each firm is represented by a
distribution of text rather than by a single characteristic.
For a target firm, the economic question is which combination of other firms
best represents its information footprint within the investable universe.
A pairwise distance can identify proximity, but it cannot determine how several
candidate firms jointly represent the target.

A firm's article embeddings define an empirical characteristic law
$\widehat C_i$; its population counterpart is $C_i$.
Quadratic Wasserstein distance aligns two article clouds by minimizing their
average squared displacement in embedding units.
That pairwise alignment supplies correspondence, after which one common set of
peer weights can be chosen across all aligned article positions.

Consider a three-firm universe with target $A$ and candidate peers $B$ and $C$.
After separately aligning the article positions of $B$ and $C$ with those of
$A$, suppose the best common convex approximation assigns coefficient $0.60$ to
the aligned positions of $B$ and $0.40$ to those of $C$ at every matched index.
In the order $(A,B,C)$, the resulting target row is
$W^\flat_{A\cdot}=(0,0.60,0.40)$.
The 60--40 weights are coordinates on aligned peer positions, not probabilities
of drawing whole articles from $B$ or $C$.

This coordinate problem differs from an unrestricted Wasserstein barycentre.
An unrestricted barycentre holds input laws and barycentre weights $a_j$ fixed
and optimizes over a new centre distribution $Q$ through an objective such as
$\inf_Q\sum_{j\ne i}a_j W_2^2(C_j,Q)$
\citep{agueh_carlier_barycenters_2011}.
The present construction instead fixes the target $C_i$ and, in the empirical
implementation, fixes its pairwise transport alignments before optimizing over
simplex peer coordinates $w_i$.
Target anchoring therefore narrows the claim from constructing a new centre to
finding coordinates for an existing firm, while retaining the barycentric
geometry.

The following definition formalizes this two-stage construction.

\begin{definition}[Target-anchored Wasserstein barycentric reconstruction]
  \label{def:w-flat}
  Write the equally sized empirical clouds as
  $\widehat C_i=M^{-1}\sum_{m=1}^M\delta_{x_{im}}$.
  Let $\mathcal A_M$ be the finite set of permutations of the $M$ support
  indices.
  For every ordered pair $i\ne j$, first fix one optimal balanced assignment
  under the implementation's predetermined deterministic selection rule,
  \begin{equation}\label{eq:w2-assignment}
    \pi_{ij}\in\operatorname*{argmin}_{\pi\in\mathcal A_M}
    \frac{1}{M}\sum_{m=1}^M
    \lVert x_{im}-x_{j,\pi(m)}\rVert_2^2,
    \qquad y_{ijm}=x_{j,\pi_{ij}(m)}.
  \end{equation}
  Conditional on these target-specific assignments, define
  \begin{equation}\label{eq:w2-barycentre}
    w_i\in\operatorname*{argmin}_{w\in\Delta_{-i}}
    \frac{1}{M}\sum_{m=1}^M
    \left\lVert x_{im}-\sum_{j\ne i}w_j y_{ijm}\right\rVert_2^2,
    \qquad
    \Delta_{-i}=\{w:w_j\geq0,\ \sum_{j\ne i}w_j=1\}.
  \end{equation}
  The output is the barycentric interaction field
  $W^\flat_{ij}=(w_i)_j$ for $j\ne i$ and $W^\flat_{ii}=0$.
\end{definition}

Equation~\eqref{eq:w2-assignment} fixes a transport alignment for each ordered
target--candidate pair, conditional on the empirical clouds and their ground
metric.
Equation~\eqref{eq:w2-barycentre} then selects the target's simplex coordinates
by minimizing reconstruction loss across all aligned candidates jointly.
Because the second objective is joint, a row of $W^\flat$ is not obtained by
applying a scalar transformation separately to each pairwise distance.

Every other eligible firm enters the candidate set, self-links are excluded,
and the assignments are target specific.
The second stage uses squared reconstruction loss and simplex-normalized
weights, while equal cloud size is part of the implementation.
These features, including the rule that selects among tied optimal assignments,
are maintained design choices rather than consequences of Wasserstein distance
alone.

Geometrically, the coefficients in $w_i$ are barycentric-type coordinates on the
aligned peer positions.
Economically, we call them approximate distributional spanning weights within
the investable universe: the objective value records how well the aligned convex
span represents the target's information footprint.
A small residual supports approximate semantic substitutability of that peer
combination for the target within the maintained representation, ground metric,
and candidate universe; the simplex constraint does not assert exact spanning.
The weights do not identify causal influence, a tradable replicating portfolio,
or a literal arbitrage relation.

Because the target determines both the alignments and the reconstruction
problem, $W^\flat$ is generally asymmetric even though pairwise Wasserstein
distance is symmetric.
A large $W^\flat_{ij}$ need not imply a large $W^\flat_{ji}$: the direction
records reconstruction relevance for firm $i$.

Figure~\ref{fig:w2-operator-construction} shows how article-level correspondence
becomes firm-level reconstruction relevance.

\begin{figure}[htbp]
  \centering
  \ppfigure{1}{w2_operator_construction}
  \caption{Target-anchored reconstruction of one row of $W^\flat$.
    For target firm $i$, separate target-to-peer transport assignments establish
    article-level correspondence.
    Holding those assignments fixed, simplex coordinates combine the aligned
    peer clouds to reconstruct the target.
  The fitted coordinates form a nonnegative unit-sum row with zero self-weight.}
  \label{fig:w2-operator-construction}
  \figalttext{Four panels read left to right. Target article points are
    separately matched to three peer clouds; three decreasing simplex weights
    combine aligned peer points into a reconstruction of the target. The same
    weights become an interaction row with zero self-weight and unit row sum.
    The diagram illustrates the construction for one target rather than an
  estimated economic effect.}
\end{figure}

Read the figure from left to right: transport establishes correspondence,
simplex reconstruction selects joint coordinates, and the fitted row becomes
the target's peer average.
The simplex constraint and leave-one-out policy imply, row by row,

\[
  W^\flat_{ij}\geq0,\qquad W^\flat_{ii}=0,\qquad
  \sum_{j} W^\flat_{ij}=1.
\]

Thus the construction satisfies \Cref{def:interaction-matrix} without a
separate row-normalization step.
It avoids a conventional kernel-bandwidth choice, not researcher choices about
the representation, candidate set, or reconstruction rule.
The formal norm implication needed for the spatial multiplier is stated in the
appendix as \Cref{thm:wflat-sar-gate}; in the main argument, the economic point
is that every empirical row has the admissible peer-average interpretation.

\subsection{Comparators}\label{sec:interaction-comparators}

The three comparators vary one margin at a time: cardinal coordinates within a
selected support, the rule that maps transport geometry into weights, and the
information object that defines a relation between firms.

The equal-active-support comparator holds fixed the active peer set selected by
$W^\flat$ and changes only the cardinal magnitudes within that set.
For each row, a peer is active when $W^\flat_{ij}>10^{-8}$, the numerical
tolerance fixed in the empirical producer.
Writing $A_i=\{j\ne i:W^\flat_{ij}>10^{-8}\}$, the
equal-active-support matrix is

\[
  W^{\mathrm{eq}}_{ij}
  =\frac{\mathbf{1}\{j\in A_i\}}{|A_i|}.
\]

It has exactly the same numerical support as $W^\flat$ but weights every active
peer equally.
Its economic question is whether fitted distributional spanning weights improve
conditional spatial fit beyond selecting the active peer set.

The RBF diffusion comparator holds fixed the empirical characteristic laws,
their pairwise quadratic Wasserstein distances, and the candidate universe, but
changes how that geometry becomes a peer row.
It replaces joint target reconstruction with a separate proximity weight
for each target--candidate pair.
For $D_{ij}=W_2(\widehat C_i,\widehat C_j)$ and bandwidth $h>0$, it sets

\[
  W^h_{ij}
  =\frac{\mathbf{1}\{j\ne i\}\exp(-D_{ij}^2/h)}
  {\sum_{k\ne i}\exp(-D_{ik}^2/h)}.
\]

The application sets $h$ to the median off-diagonal value of $D_{ij}^2$.
The economic question is whether peer dependence is organized by pairwise
proximity alone or by the joint reconstruction of a fixed target.
This dense operator also makes the conventional bandwidth choice explicit, and
the appendix documents its construction.

The persistent news co-mention comparator holds fixed the firm universe, the
2018--2022 pre-evaluation timing, and the row-stochastic econometric role, but
changes the information object and the link rule.
The matrix $W^{\mathrm{news}}$ links firms that are the only two tagged names in
an article and co-occur in at least two calendar years during that period.
Total co-mention counts are symmetrized and row normalized.
Its economic question is whether interaction is organized by approximate
distributional spanning or by persistent shared news coverage.
Because it is admissible in the sense of \Cref{def:interaction-matrix}, it can
also enter \Cref{def:two-field-adjustment} as a separate adjustment field rather
than only as an alternative to $W^\flat$.
\Cref{sec:results} uses it both ways: first as a comparator estimated on its
own, then as the second channel of the joint model, where the two roles are the
boundary and the interior of one nested family.

The barycentric field establishes empirical peers from observable
characteristic laws.
It does not yet determine what separation between those laws implies for latent
risk exposures.
The next section states that characteristic-to-exposure restriction and asks how
much of the implied exposure dispersion survives equilibrium peer adjustment.
